% Part: first-order-logic
% Chapter: introduction
% Section: satisfaction

\documentclass[../../../include/open-logic-section]{subfiles}

\begin{document}

\olfileid{fol}{int}{sat}

\olsection{సంతృప్తి}

\tetoken{సూత్రాలు}{!!{formula}s} అంటే ఏమిటో తెలిసిన వెంటనే మొదటిస్థాయి విధేయ తర్కపు
అర్థవిచారానికి ముందుకు వెళ్లవచ్చు: ఇక్కడ ప్రాథమిక నిర్వచనం
\tetoken{ఒక నిర్మాణానికి}{!!a{structure}} సంబంధించినది. మన సరళ భాషకు ఒక \tetoken{నిర్మాణం}{!!a{structure}}~$\Struct M$లో
కేవలం మూడు అంశాలు ఉంటాయి: \emph{\tetoken{వ్యక్తి క్షేత్రం}{!!{domain}}} అనే ఖాళీ కాని
సమితి~$\Domain{M}$, $\Struct M$లో $\Obj a$ ఎంచి చూపే వస్తువు,
$\Struct M$లో $\Obj P$ వర్తించే వస్తువులు. $\Obj a$ ఎంచి చూపే
వస్తువును $\Assign{\Obj a}{M}$తో, $\Obj P$ వర్తించే వస్తువుల
సమితిని $\Assign{\Obj P}{M}$తో సూచిస్తాం. \tetoken{ఒక నిర్మాణం}{!!^a{structure}}~$\Struct{M}$
ఈ మూడు విషయాలతోనే ఏర్పడుతుంది: $\Domain{M}$,
$\Assign{\Obj a}{M} \in \Domain{M}$, $\Assign{\Obj P}{M}
\subseteq \Domain{M}$. సాధారణ సందర్భం మరింత సంక్లిష్టంగా ఉంటుంది;
ఎందుకంటే అనేక \tetoken{విధేయ సంకేతాలు}{!!{constant}s}, \tetoken{స్థిరసంకేతాలు}{!!{constant}s} ఉంటాయి,
\tetoken{విధేయ సంకేతాలకు}{!!{predicate}s} ఒకటి కంటే ఎక్కువ స్థానాలు ఉండవచ్చు, అలాగే
\tetoken{ప్రమేయ సంకేతాలు}{!!{function}s} కూడా ఉంటాయి.
\sourcecorrection{OLTEINT-003}{ఆంగ్ల మూలం బహుస్థానాలుగా ఉండగలవని
స్థిరాంకాల గురించి చెప్పింది. వెంటనే ముందు, తరువాత ఇచ్చిన నిర్మాణ
వివరణల ప్రకారం బహుస్థాన అరిటీ విధేయాలకు చెందుతుంది; స్థిరాంకాలకు
కాదు. అందువల్ల ఆ ఒక్క వర్గనామాన్ని విధేయాలుగా సరిచేశాం.}

\tetoken{చరాలు}{!!{variable}s} లేని \tetoken{సూత్రాలకు}{!!{formula}s} సంతృప్తి నిర్వచనం ఇవ్వడానికి ఇది
చాలు. \tetoken{సూత్రాలను}{!!{formula}s} మనం నిర్వచించిన విధానాన్నే ప్రతిబింబించే
ఆగమనాత్మక నిర్వచనం ఇవ్వాలన్నది భావన. ఒక పరమాణు సూత్రం
  $\Struct{M}$లో ఎప్పుడు సంతృప్తి చెందుతుందో పేర్కొని, తరువాత ఉదాహరణకు
  $\Struct{M}$లో $!A$ సంతృప్తి చెందుతుందా లేదా అన్నదాని ఆధారంగా $\lnot !A$ ఎప్పుడు
$\Struct{M}$లో సంతృప్తి చెందుతుందో చెబుతాం. ఉదాహరణకు, ఇలా
నిర్వచించవచ్చు:
\begin{enumerate}
  \item $\Atom{\Obj P}{\Obj a}$, $\Struct{M}$లో సంతృప్తి చెందడం,
  $\Assign{\Obj a}{M} \in \Assign{\Obj P}{M}$ అయినప్పుడూ, అప్పుడు
  మాత్రమే.
  \item $\lnot !A$, $\Struct{M}$లో సంతృప్తి చెందడం, $\Struct{M}$లో $!A$
  సంతృప్తి చెందనప్పుడూ, అప్పుడు మాత్రమే.
  \item $(!A \land !B)$, $\Struct{M}$లో సంతృప్తి చెందడం,
  $\Struct{M}$లో $!A$, అలాగే $\Struct{M}$లో $!B$ కూడా సంతృప్తి
  చెందినప్పుడూ, అప్పుడు మాత్రమే.
\end{enumerate}
$\Domain{M} = \{0, 1, 2\}$, $\Assign{\Obj a}{M} = 1$,
$\Assign{\Obj P}{M} = \{1, 2\}$ అని అనుకుందాం. ఈ నిర్వచనం
$\Atom{\Obj P}{\Obj a}$, $\Struct{M}$లో సంతృప్తి చెందుతుందని
చెబుతుంది (ఎందుకంటే $\Assign{\Obj a}{M} =  1 \in \{1,2\} =
\Assign{\Obj P}{M}$). ఇంకా $\lnot \Atom{\Obj P}{\Obj a}$,
$\Struct{M}$లో సంతృప్తి చెందదని; దాని వల్ల $\lnot\lnot
\Atom{\Obj P}{\Obj a}$ సంతృప్తి చెందుతుందని, $(\lnot
\Atom{\Obj P}{\Obj a} \land \Atom{\Obj P}{\Obj a})$ సంతృప్తి
చెందదని; ఇలాగే మరిన్ని విషయాలను చెబుతుంది.

పరిమాణీకరణికాలకు నిర్వచనం ఇవ్వాలనుకున్నప్పుడు ఇబ్బంది వస్తుంది:
``$\lexists[\Obj v_0][\Atom{\Obj P}{\Obj v_0}]$ సంతృప్తి చెందడం,
$\Atom{\Obj P}{\Obj v_0}$ సంతృప్తి చెందినప్పుడూ, అప్పుడు మాత్రమే''
అని చెప్పాలని అనిపిస్తుంది. కానీ \tetoken{చరాలను}{!!{variable}s} ఎలా పరిగణించాలో
\tetoken{నిర్మాణం}{!!{structure}}~$\Struct{M}$ మనకు చెప్పదు. నిజానికి మనం చెప్పాలనుకున్నది,
$\Atom{\Obj P}{\Obj v_0}$ అనేది \emph{$\Obj v_0$ యొక్క ఏదో ఒక
విలువకు} సంతృప్తి చెందుతుందని. దీన్ని ఖచ్చితపరచడానికి,
$\Domain{M}$లోని \tetoken{మూలకాలను}{!!{element}s} $\Obj a$కే కాకుండా $\Obj v_0$కు
కూడా నిర్దేశించే పద్ధతి కావాలి. అందుకోసం \emph{\tetoken{చరం}{!!{variable}}
నిర్దేశాలను} ప్రవేశపెడతాం. \tetoken{ఒక చరం}{!!^a{variable}} నిర్దేశం అంటే,
\tetoken{చరాలను}{!!{variable}s} $\Domain{M}$లోని \tetoken{మూలకాలకు}{!!{element}s} పంపే ఒక ప్రమేయం~$s$
(మన ఉదాహరణలో $0$, $1$, $2$లలో ఒకదానికి పంపుతుంది). ఒక
\tetoken{సూత్రంలో}{!!a{formula}} ఏ \tetoken{చరాలు}{!!{variable}s} రావచ్చో ముందుగా తెలియదు కాబట్టి, $s$
ఏ \tetoken{చరాలకు}{!!{variable}s} విలువలు నిర్దేశిస్తుందో పరిమితం చేయలేం. $s$ అన్ని
\tetoken{చరాలు}{!!{variable}s}~$\Obj v_0$, $\Obj v_1$, \dotsకు విలువలు నిర్దేశించాలని
కోరడం సరళమైన పరిష్కారం. అవసరమైన వాటినే మనం ఉపయోగిస్తాం.
\sourcecorrection{OLTEINT-004}{ఆంగ్ల మూలం $\Domain{M}=\{0,1,2\}$గా
నిర్ణయించిన తరువాత, చరరాశి నిర్దేశం $1$, $2$, $3$లలో ఒకదానికి
పంపుతుందని చెప్పింది. నిర్మాణపు అదే ప్రకటిత ప్రదేశానికి అనుగుణంగా
జాబితాను $0$, $1$, $2$గా సరిచేశాం.}

\tetoken{సూత్రాలు}{!!{formula}s} సంతృప్తిని కేవలం \tetoken{ఒక నిర్మాణానికి}{!!a{structure}} సాపేక్షంగా
నిర్వచించకుండా, \tetoken{ఒక నిర్మాణం}{!!a{structure}}~$\Struct{M}$కు \emph{మరియు}
\tetoken{ఒక చరం}{!!a{variable}} నిర్దేశం~$s$కు సాపేక్షంగా నిర్వచించి, సంక్షిప్తంగా
$\Sat{M}{!A}[s]$ అని రాస్తాం. \tetoken{చరాలు}{!!{variable}s} ఉన్న పరమాణు
\tetoken{సూత్రాలను}{!!{formula}s} నిర్వహించడానికి మన నిర్వచనంలో ఇప్పుడు మరో ఉపవాక్యం
ఉంటుంది:
\begin{enumerate}
  \item $\Sat{M}{\Atom{\Obj P}{\Obj a}}[s]$ కావడం,
  $\Assign{\Obj a}{M} \in \Assign{\Obj P}{M}$ అయినప్పుడూ, అప్పుడు
  మాత్రమే.
  \item $\Sat{M}{\Atom{\Obj P}{\Obj v_i}}[s]$ కావడం,
  $s(\Obj v_i) \in \Assign{\Obj P}{M}$ అయినప్పుడూ, అప్పుడు మాత్రమే.
  \item $\Sat{M}{\lnot !A}[s]$ కావడం, $\Sat{M}{!A}[s]$ కాకపోయినప్పుడూ,
  అప్పుడు మాత్రమే.
  \item $\Sat{M}{(!A \land !B)}[s]$ కావడం, $\Sat{M}{!A}[s]$ మరియు
  $\Sat{M}{!B}[s]$ అయినప్పుడూ, అప్పుడు మాత్రమే.
\end{enumerate}
సరే, ఇది ఒక సమస్యను పరిష్కరిస్తుంది: $s(\Obj v_0)$ అనే విలువకు
$\Struct{M}$, $\Atom{\Obj P}{\Obj v_0}$ను ఎప్పుడు సంతృప్తిపరుస్తుందో
ఇప్పుడు చెప్పవచ్చు. $\lexists[\Obj v_0][\Atom{\Obj P}{\Obj v_0}]$కు
నిర్వచనాన్ని సరిగ్గా పొందడానికి ఇంకొక పని చేయాలి:
$\Sat{M}{\lexists[\Obj v_0][\Atom{\Obj P}{\Obj v_0}]}[s]$ కావడం,
$\Obj v_0$కు విలువను నిర్దేశించే \emph{ఏదో ఒక} పద్ధతి~$s'$కు
$\Sat{M}{\Atom{\Obj P}{\Obj v_0}}[s']$ అయినప్పుడూ, అప్పుడు మాత్రమే
కావాలని మనం కోరుతున్నాం. అయితే $\Obj v_0$కు నిర్దేశించిన విలువ,
$s(\Obj v_0)$ ఎంచి చూపే విలువే కావాల్సిన అవసరం లేదు. దానికి ఒక
సంజ్ఞావిధానాన్ని ప్రవేశపెడతాం: $m \in \Domain{M}$ అయితే,
$\Subst{s}{m}{\Obj v_0}$ అనేది $s$లాగే ఉండే నిర్దేశం (అంటే $\Obj
v_0$ కాకుండా మిగతా అన్ని \tetoken{చరాలపై}{!!{variable}s} అలాగే ఉంటుంది), కానీ
$\Obj v_0$కు మాత్రం $m$ను నిర్దేశిస్తుంది. ఇప్పుడు మన నిర్వచనం ఇలా
ఉండవచ్చు:
\begin{enumerate}\setcounter{enumi}{4}
  \item $\Sat{M}{\lexists[\Obj v_i][!A]}[s]$ కావడం,
  ఏదో~$m \in \Domain{M}$కు $\Sat{M}{!A}[\Subst{s}{m}{\Obj v_i}]$
  అయినప్పుడూ, అప్పుడు మాత్రమే.
\end{enumerate}
ఇది పనిచేస్తుందా? ప్రతి~$i \in \Nat$కు $s(\Obj v_i) = 0$గా
తీసుకుందాం. $\Sat{M}{\lexists[\Obj v_0][\Atom{\Obj P}{\Obj
v_0}]}[s]$ కావడం, $\Sat{M}{\Atom{\Obj P}{\Obj
v_0}}[\Subst{s}{m}{\Obj v_0}]$ అయ్యే $m \in \Domain{M}$ ఏదైనా
ఉన్నప్పుడూ, అప్పుడు మాత్రమే. అలాంటి విలువ ఉంది: $m = 1$ లేదా $m = 2$ను
ఎంచుకోవచ్చు. $s$ స్వయంగా $\Obj v_0$కు నిర్దేశించిన
$s(\Obj v_0)$ విలువ---ఈ సందర్భంలో $0$---ఆ పని చేయకపోయినా ఇది
సత్యమేనని గమనించండి. మనకు $\Sat{M}{\Atom{\Obj P}{\Obj
v_0}}[\Subst{s}{1}{\Obj v_0}]$ ఉంటుంది; కానీ
$\Sat{M}{\Atom{\Obj P}{\Obj v_0}}[s]$ ఉండదు.

ఇది గందరగోళంగా, భారంగా కనిపిస్తే: నిజంగానే అలాగే ఉంటుంది. కానీ అన్ని
\tetoken{సూత్రాలకు}{!!{formula}s} సంతృప్తిని ఖచ్చితంగా, ఆగమనాత్మకంగా నిర్వచించడానికి ఈ
అదనపు సంక్లిష్టత అవసరం; అర్థపర భావనలను ఖచ్చితంగా నిర్వచించడానికి
మనకు ఇలాంటిదొకటి కావాలి. ఇతర పద్ధతులు కూడా ఉన్నాయి; కానీ అవన్నీ
సమానంగా (అ)సొగసైనవే.

\end{document}
